\section{The order-336 source and effective the quadratic field of discriminant minus seven bounds}
\label{sec:g24-effective}

\paragraph{HD-S13 (exact group, primitive projector, and integral order).}
Put
\[
 K=\mathbb Q(\lambda),\qquad \lambda^2+\lambda+2=0,\qquad
 \lambda=\frac{-1+\sqrt{-7}}2,\qquad\mathcal O=\mathbb Z[\lambda].
\]
Let \(\mathcal G\) be generated by the three matrices in (G12.1), and
write \(G=336\). These matrices preserve
\[
 H=\begin{pmatrix}2&-1&-\lambda\\-1&2&-1\\1+\lambda&-1&2\end{pmatrix}.
\]
Its leading principal minors are \(2,3,1\); it is positive definite.
The exact integer-ring certificate closes the generated group after
336 distinct matrices and verifies inverses. Its element-order counts
are
\[
\begin{array}{c|rrrrrrr}
\text{order}&1&2&3&4&6&7&14\\ \hline
\text{number}&1&43&56&84&56&48&48.
\end{array}
\]
The reflection character \(\chi\) has degree three,
\(\sum_{g\in\mathcal G}|\chi(g)|^2=336\), and a value \(\lambda\).
It is therefore absolutely irreducible and has character field
exactly \(K\); the displayed matrices realize it over \(K\).

The connection to published root data is an explicit conjugation.
Cohen, \emph{Finite complex reflection groups},
Ann.\ Sci.\ ENS (4) 9 (1976), (4.15), printed p.407, gives unit
roots \(e_1,e_2,e_3\) with his parameter
\(\alpha=1+\lambda\). Set
\[
 2T=\begin{pmatrix}0&-1&0\\2&-1&-\lambda\\
                   0&1+\lambda&-\lambda\end{pmatrix}.
\]
The columns of \(T\) are \(e_1,e_3,e_2\), and direct multiplication
gives \(T^*T=H/2\). Let \(u_j\) denote the standard coordinate
vectors in \(K^3\). Then
\(Ts_jT^{-1}=1-2(Tu_j)(Tu_j)^*\), the corresponding Cohen
unit-root reflections. Cohen's (4.16), table on p.412, identifies
this group as Shephard--Todd \(G_{24}\), with degrees \(4,6,14\)
and 21 reflections. The enumeration and subsequent proofs use
the displayed exact matrices themselves.

Let \(C\) be a cover from HD-G12 and let \(J=J(C)\).
Set \(h=s_1s_2\). The exact matrices give \(h^3=1\),
\(\chi(h)=\chi(h^2)=0\), so the averaging projector
\((1+h+h^2)/3\) has rank one in the reflection representation.
The rational central idempotent of the conjugate character pair is
\[
 \varepsilon=\frac3G\sum_{g\in\mathcal G}
       \operatorname{Tr}_{K/\mathbb Q}\chi(g^{-1})\,g.
\]
The representation is realized over \(K\), so its rational simple
factor is \(M_3(K)\). The element
\[
 e=\varepsilon\,\frac{1+h+h^2}{3}                           \tag{S13.1}
\]
is a primitive idempotent in this factor. It is fixed by the
Rosati involution \(g\mapsto g^{-1}\), since central projection and
subgroup averaging are self-adjoint and commute.
The factor three in \(\varepsilon\) cancels the averaging
denominator: the coefficient of \(g\) in \(Ge\) is
\[
 \sum_{j=0}^2\operatorname{Tr}_{K/\mathbb Q}
                       \chi(h^j g^{-1})\in\mathbb Z.       \tag{S13.2}
\]
Thus \(Ge\in\mathbb Z[\mathcal G]\).

Define the actual abelian subvariety
\(A_C=\operatorname{Im}(Ge:J\to J)^0\), with
\(\theta_C=\Theta_J|_{A_C}\).
The Chevalley--Weil calculation in HD-G12 gives multiplicities
three in \(H^{1,0}\) for each conjugate character; the rank-one
matrix idempotent selects a Weil sixfold of signature \((3,3)\).
Its integral lattice is
\(\Lambda_C=H_1(J,\mathbb Z)\cap eH_1(J,\mathbb Q)\).
Self-adjointness, the primitive lattice inclusion, and
unimodularity of \(H_1(J,\mathbb Z)\) give
\[
 \Lambda_C^\#=eH_1(J,\mathbb Z),\qquad
 G\Lambda_C^\#\subseteq\Lambda_C.
\]
Indeed, any integral functional on \(\Lambda_C\) extends to the
unimodular ambient lattice, and its orthogonal projection by \(e\)
represents it on \(\Lambda_C\). Every polarization elementary
divisor therefore divides \(G\). In particular,
\[
 \Delta_C\mid G^6,\qquad \int_{A_C}\theta_C^6\le6!G^6.       \tag{S13.3}
\]

For the integral field action, take \(v=(s_1s_2s_3)^2\).
It has order seven and \(\chi(v)=\lambda\).
Its conjugacy class has 24 elements, and the exact class sum acts
as \(8\lambda\,1_3\) in the reflection representation. This also
follows from Schur's lemma and \(24\lambda/3=8\lambda\).
The class sum belongs to \(\mathbb Z[\mathcal G]\), is central,
and preserves \(A_C\). Thus \(8\lambda\) acts integrally, giving
the actual order
\[
 \mathcal O_8=\mathbb Z+8\mathbb Z[\lambda]
 \subseteq\operatorname{End}(A_C).                        \tag{S13.4}
\]
The Rosati involution sends this class sum to its inverse-class
sum and restricts to conjugation on \(K\).

\paragraph{HD-S14 (a uniform explicit detector on every split target).}
For the rest of this section the following integers are fixed:
\[
\begin{gathered}
 G=336,\qquad g=2G+1=673,\qquad
 R=(2G+1)G!G^G,\qquad L=2gR,\\
 S_0=G^{2g}(g!)^2,\qquad S_B=(L/2)!S_0^R,\qquad k=L/2+3,\\
 B_6=(24G)^k\binom{k}{3}S_B\,6!G^6,\\
 \boxed{\mathcal D_7(q)=
       8^3\bigl(4(8G)^6q\bigr)^6 B_6.}                   \tag{S14.1}
\end{gathered}
\]
These are explicit positive integers, with factorials left in exact
form. For every polarized split Weil sixfold with integral
\(\mathcal O\)-action and polarization product \(\Delta\),
\[
 \boxed{\delta_W\le\mathcal D_7(\Delta).}                  \tag{S14.2}
\]
For a fixed order of conductor \(f\), the corresponding statement is
\[
 \boxed{\delta_W\le f^{33}\mathcal D_7(\Delta).}            \tag{S14.3}
\]
Here split means ordinary Hermitian determinant class \([-1]\);
HD-S03a proves the equivalence in signature \((3,3)\).

\paragraph{Proof: all supplier cycles and the source degree.}
The inertia of \(C\to\mathbb P^1\) consists of twelve reflections
of order two, so \(g(C)=2G+1\). For every linear character
\(\psi\) of a subgroup \(H'\subset\mathcal G\), the actual cyclic
cover
\[
 C/\ker\psi\longrightarrow C/H'
\]
has inertia of order one or two. A nontrivial faithful branch
exponent is consequently \(m/2\), where \(m=\operatorname{ord}\psi\).
The branch number is even by the cyclic product relation. Thus
every retained nonreal supplier has positive even multiplicity
\(a=2g(C/H')-2+r\). Its branch exponents are paired with their
negatives.

The proved paired-involution supplier argument HD-G05 applies
word for word to these actual quotient covers. Its ingredients
are Schoen1988, Theorem 2.0 and Lemmas 2.3--2.6, printed
pp.11--17, and the explicitly proved exceptional-divisor vanishing:
on every nonvertical exceptional \(\mathbb P^1\)-fibre the faithful
eigensheaf restricts to a rank-one local system on
\(\mathbb C^\times\) with monodromy \(-1\), whose cohomology is
zero; on a vertical branch fibre the faithful eigensheaf is zero.
Proper base change makes the exceptional summand vanish. The
unramified case is Schoen's original \(r=0\) construction.
This gives actual nonzero determinant cycle components for every
supplier required below.

Brauer induction, followed by the determinant-character
homomorphism for the full symplectic centralizer \(U\) of
\(\mathcal G\), expresses the target determinant character as
an integral combination of these supplier characters.
The full-centralizer proof HD-G06 includes the order-two
orthogonal character coordinates. Terms of order one or two
have zero determinant coordinate: conjugate complex
multiplicities agree and quaternionic multiplicities are even.
Zero-multiplicity terms also contribute zero. The retained
suppliers are precisely those treated above.

For completeness, the coefficient bound does not depend on a
deposited character identity. Append columns \(2e_j\) for the
order-two coordinates. There are at most \(G\) rows, column
entries have absolute value at most \(G\), and the target has
coordinate size at most one. For an independent square block
of rank \(r\), its determinant has absolute value at most
\(r!G^r\). A shortest path in the quotient lattice has fewer
than that many steps; Cramer's rule bounds the remaining
coefficients by \(r!G^r\). The total absolute coefficient is at
most \((r+1)r!G^r\le R\). This is the explicit integral
linear-algebra proof already given in SQ3 and HD-G07.

For a supplier of cyclic order \(m\) and multiplicity \(a\),
Schoen's quotient \(C'^a/G_0\), with
\(|G_0|=m^{a-1}a!\), contains a degree-one lift of a linear
space \(\mathbb P^{a/2}\). Pulling an irreducible component
to \(C'^a\), taking the complete orbit sum, and applying the
integral Ramanujan numerator for the primitive characters
gives an integral signed cycle of point-polarization mass
at most
\[
 m^{3a/2}(a!)^2\le S_0.                                  \tag{S14.4}
\]
Here \(a\le g(C')\le g\) and \(m\le G\).
This bound follows from the actual quotient degrees:
the pulled incidence divisor is \(m\) times the sum of point
classes, the chosen component has degree at most \(|G_0|\)
over the lifted linear space, the orbit sum costs another
\(|G_0|\), and the Ramanujan coefficients have total absolute
value at most \(m^2\). It uses no flatness assertion along
the lifted linear space.

Exterior products over the bounded Brauer identity use an even
number \(\ell\le gR\le L\) of curve factors, with total mass
at most \(S_B\). Choose the conjugate determinant line for each
negative coefficient. Poincare duality on these middle curve-power
tensors gives a \(U\)-invariant correspondence whose action sends
the supplier product to a nonzero vector in either target Weil
eigenline. The explicit tensor-and-contraction proof is HD-G06
and SQ5; all degrees here are \(\ell+6\), so no inverse
hard-Lefschetz map occurs.

Milne1999, Lemma 3.5, Proposition 3.6 and Remark 3.7, applied
to the full centralizer \(U\), show that this invariant
correspondence is a complex linear combination of products
of degree-two invariant classes. The proof in HD-G06 supplies
the actual generating divisor maps on \(J(C)^{\ell+1}\):
\[
 \operatorname{pr}_j^*(2\Theta_J),\qquad
 (\operatorname{pr}_j+a\operatorname{pr}_{j'})^*(2\Theta_J),
 \quad a\in\mathcal G,
\]
including \(j=j'\). These line bundles are globally generated.
The coefficient sizes of their complex spanning combination
are unnecessary: one of the resulting actual effective cycles
must have nonzero Weil projection.

Here are the precise numerical replacements for HD-G07.
For a quotient map \(C\to C'\) of degree \(b\), the
Abel--Jacobi map followed by pullback to \(J(C)\) has
theta degree
\[
 b\,g(C')\le g(C)-1+b\le3G.
\]
The theorem of the square gives a factor eight for the
displayed divisor systems, hence on the product of the curve
supplier and \(A_C\) their classes are bounded in the nef order
by \(24G(H_B+\theta_C)\).
Choose successive general members so that all intersections
with the finitely many effective components are proper.
For \(c=\ell/2+3\le k\), each resulting pushforward has mass
at most
\[
 (24G)^c\binom c3 S_B\int_{A_C}\theta_C^6
 \le B_6.
\]
The binomial coefficient selects the unique surviving mixed
degree. Nonvanishing of the original invariant contraction
proves that some reduced irreducible effective three-cycle
detects the Weil plane. This constructs a detector of degree
at most \(B_6\) on every source \(A_C\).

\paragraph{Proof: the two actual isogenies and all fibres.}
The source order in (S13.4) is sufficient to construct its
maximal-order saturation:
\[
 \Lambda_C'=\mathcal O\Lambda_C
           =\Lambda_C+\lambda\Lambda_C,\qquad
 q:A_C\longrightarrow A_C'.
\]
One has \(8\Lambda_C'\subseteq\Lambda_C\).
The form \(8E_C\) is integral on \(\Lambda_C'\):
the mixed terms use the integral \(8\lambda\) action, and
\(E_C(\lambda x,\lambda y)=2E_C(x,y)\).
It remains a positive Riemann form of type \((1,1)\).
If \(d=\deg q\), its polarization satisfies
\[
 q^*\theta_C'=8\theta_C,\qquad
 \Delta_C'=\frac{8^6\Delta_C}{d}\mid(8G)^6.                \tag{S14.5}
\]
Pushforward of a detecting three-cycle is effective, preserves
its nonzero Weil projection under this \(K\)-linear isogeny,
and multiplies its intersection degree by at most \(8^3\).

The ring \(\mathcal O\) is a PID. This can be seen directly
from its norm-Euclidean algorithm: for a complex number
\(x+y\lambda\), choose \(b\in\mathbb Z\) nearest to \(y\),
then \(a\in\mathbb Z\) nearest to \(x-(y-b)/2\).
The remainder has norm
\((x-a-(y-b)/2)^2+7(y-b)^2/4\le
1/4+7/16=11/16<1\).
For a split integral Hermitian lattice with polarization
product \(q\), a primitive isotropic half-lattice extends
to a basis, with Gram matrix
\(\begin{psmallmatrix}0&B\\B^*&D\end{psmallmatrix}\) and
\(N(\det B)=q\). The substitutions
\[
 w'=w-v(B^*)^{-1}D/2,\qquad u=w'B^{-1}
\]
give a hyperbolic basis. Both this coordinate map and its
inverse are cleared by \(2q\). The identity Hermitian matrix
gives a unimodular alternating form for the integral basis
\((v_j,\lambda v_j)\), so the polarization normalization is
again \(|\det H|=q\), with no hidden field-discriminant factor.

HD-G12 proves that the rational Hermitian space of \(A_C'\)
is split and that its period map has rank nine.
Comparing (S14.5) to any fixed target lattice of product
\(\Delta\), the explicit denominator
\[
 N=4(8G)^6\Delta
\]
clears both rational coordinate maps. It gives actual
isogenies with
\(\phi^*\theta_{\mathrm{target}}=N^2\theta_C'\).
Pushforward of a three-cycle costs \(N^6\), giving exactly
\(\mathcal D_7(\Delta)\). The rank-nine image contains an
analytic open subset of the fixed polarized Weil moduli
component. The supplied bounded-detector closedness theorem
extends this same bound to every fibre of that component,
proving (S14.2).

For a target with integral order
\(\mathcal O_f=\mathbb Z+f\mathbb Z[\lambda]\), saturate
\(\Lambda\) to \(\Lambda'=\Lambda+\lambda\Lambda\) and
use \(fE\). The mixed terms are integral and the
\(\lambda,\lambda\) term is \(2fE(x,y)\), so the resulting
quotient \(q_f\) is actually polarized, with
\[
 \Delta'=\frac{f^6\Delta}{d},\qquad
 \deg_\theta q_f^*Z'=\frac d{f^3}\deg_{\theta'}Z'.
\]
The function \(\mathcal D_7\) is homogeneous of degree six.
Thus the exact bound is
\[
 \frac d{f^3}\mathcal D_7\left(\frac{f^6\Delta}{d}\right)
 =\frac{f^{33}}{d^5}\mathcal D_7(\Delta)
 \le f^{33}\mathcal D_7(\Delta),
\]
proving (S14.3). \(\square\)

\paragraph{HD-S15 (small integral extraction for this field).}
On a maximal-order sixfold put \(A_\lambda=[\lambda]^*\),
\(\bar A_\lambda=[\bar\lambda]^*\) and \(S=A_\lambda+\bar A_\lambda\).
On the seven embedding summands of \(H^6\), the eigenvalues of
\(S\), paired by conjugation, are \(9,2,-12,16\).
Thus
\[
 \begin{split}
 P&=-(S-2)(S+12)(S-16)=1029e_W\\
  &=384-8(A_\lambda+\bar A_\lambda)
       +6(A_\lambda^2+\bar A_\lambda^2)
       -(A_\lambda^3+\bar A_\lambda^3).
 \end{split}
\]
The second identity uses
\(A_\lambda\bar A_\lambda=[2]^*=64\) on \(H^6\).
It specifies an actual integral cycle operator; no equality
on Chow groups is asserted. Its two signed parts each have
degree at most
\(384+12\cdot64=16\cdot8+2\cdot512=1152\)
times the detector degree.
Since \(\lambda^6=2-5\lambda\), the rotation
\(Z=2A_\lambda-9\) is skew-adjoint on \(W\), with
\(Z^2=-175\) and degree cost \(L=25\).
The exact Gram extraction gives
\[
 A_0=5147144,\qquad
 A_1=\left\lceil19756800\sqrt{14}\right\rceil=73923177.
\]
The exact integer-pair script checks every eigenvalue and
both strict/weak inequalities defining this ceiling.
Replacing each endomorphism term \([u]^*\) by \([fu]^*\)
gives the explicit integral-order factor \(f^{24}\), as in
HD-D03. Therefore
\[
 \boxed{c_t\le73923177\,\delta M_T f^{90}
                  \mathcal D_7(\Delta)^2}                \tag{S15.1}
\]
for every split sixfold of this field and conductor \(f\).
Here \(M_T=\lceil\sqrt{-\int t^2}\rceil\), and
\(90=24+2\cdot33\).

\paragraph{HD-S16 (all determinant fourfolds over the same field).}
Let \(A\) be a Weil fourfold for \(K\), with maximal-order
action, polarization product \(\Delta\), and arbitrary ordinary
determinant class. Append the actual balanced surface
\(B=E_K^2\), with action \((\lambda,\bar\lambda)\) and
polarization \(\eta_1+\Delta\eta_2\).
The compatible Hermitian matrices have determinants
\(\Delta\) and \(-\Delta\), so \(A\times B\) is a split
sixfold of polarization product \(\Delta^2\).
It has a detector of degree \(\mathcal D_7(\Delta^2)\)
by (S14.2).

Write \(\alpha=1+\lambda\), and use the four actual row maps
\(x,y,x+y,x+\alpha y\) on \(B\).
Their pullback point divisors have intersection matrix
\[
 \begin{pmatrix}
 0&1&1&2\\1&0&1&1\\1&1&0&2\\2&1&2&0
 \end{pmatrix},\qquad\det=-7.
\]
They span the full rational \((1,1)\) space as before.
Twice these classes are globally generated and at most
\(6\theta_B\) in the nef order.
For the fourth class, the relevant difference matrix is
\[
 6\begin{pmatrix}1&0\\0&\Delta\end{pmatrix}
 -2\begin{pmatrix}1&\alpha\\\bar\alpha&2\end{pmatrix},
\]
with nonnegative principal minors \(4,6\Delta-4\) and
\(24(\Delta-1)\). The other three classes satisfy the
smaller bound \(4\theta_B\).

The full Kunneth projection and four separating contractions
in HD-S11 now apply with these specified maps and pairings.
A proper globally generated divisor intersection and
pushforward gives an effective detecting surface of degree
at most
\[
 6\int_Yp_A^*\theta_A^2p_B^*\theta_B
 \le2\int_Y(p_A^*\theta_A+p_B^*\theta_B)^3
 \le2\mathcal D_7(\Delta^2).
\]
This proves
\[
 \boxed{\delta_W(A)\le2\mathcal D_7(\Delta^2).}             \tag{S16.1}
\]
Saturating an order of conductor \(f\) and pulling back
surfaces multiplies this bound by at most \(f^{46}\):
the exact factor is \(f^{46}/d^{11}\), since
\(\mathcal D_7(\Delta^2)\) has degree twelve in \(\Delta\).
These are all determinant fourfold bounds, obtained by the
actual product and contraction maps displayed above.

The corresponding small fourfold projector is equally explicit.
On \(H^4\), the operator \(S=A_\lambda+\bar A_\lambda\) has
eigenvalues \(1,-6,8\), paired by conjugation. Hence
\[
 49e_W=-(S+6)(S-8)
      =16+2(A_\lambda+\bar A_\lambda)
           -(A_\lambda^2+\bar A_\lambda^2).
\]
Its signed parts each have degree multiplier \(32\).
The rotation \(2A_\lambda-1\) has square \(-63\) on \(W\)
and degree cost \(9\), since \(\lambda^4=2+3\lambda\).
Thus the exact extraction coefficients are
\[
 A_0=37816,\qquad
 A_1=\lceil28224\sqrt7\rceil=74674.
\]
The integer-pair certificate checks all five eigenvalues and
the exact ceiling. Scaling the actual endomorphism terms by
\(f\) gives the factor \(f^{16}\). With
\(M_T=\lceil\sqrt{\int_A t^2}\rceil\), this proves
\[
 \boxed{c_t\le298696\,\delta M_T f^{108}
                   \mathcal D_7(\Delta^2)^2.}             \tag{S16.2}
\]

